\documentclass[12pt]{article}

\usepackage[margin=1in]{geometry}
\usepackage{setspace}
\usepackage{amsmath,amssymb,amsthm}
\usepackage{graphicx}
\usepackage{booktabs}
\usepackage{array}
\usepackage{multirow}
\usepackage{threeparttable}
\usepackage{natbib}
\usepackage{hyperref}
\usepackage[table]{xcolor}
\usepackage{caption}
\usepackage{float}
\usepackage{placeins}
\usepackage{enumitem}
\usepackage{titlesec}
\usepackage{siunitx}

\doublespacing
\hypersetup{colorlinks=true,linkcolor=black,citecolor=black,urlcolor=blue}
\captionsetup{font=small,labelfont=bf,justification=raggedright,singlelinecheck=false}
\setlength{\parindent}{0.5in}
\setlength{\parskip}{0pt}

\titleformat{\section}[block]{\normalfont\normalsize\bfseries\centering}
  {\Roman{section}.}{0.5em}{\MakeUppercase}
\titleformat{\subsection}[block]{\normalfont\normalsize\itshape}
  {\Alph{subsection}.}{0.5em}{}
\titleformat{\subsubsection}[runin]{\normalfont\normalsize\itshape}
  {}{0em}{}[.---]
\titlespacing{\section}{0pt}{18pt}{6pt}
\titlespacing{\subsection}{0pt}{12pt}{4pt}

\definecolor{darkblue}{RGB}{10,30,80}

% ─────────────────────────────────────────────────────────────────────────────
\begin{document}

% ── Title Page ────────────────────────────────────────────────────────────────
\begin{titlepage}\begin{center}\vspace*{1.2cm}

{\Large\bfseries Fire Without Smoke:\\[8pt]
Geopolitical Risk and the New Inflation Decoupling}

\vspace{1.4cm}
{\normalsize [Author Name]%
\footnote{[Department], [University].
E-mail: \texttt{author@university.edu}.
Replication materials: \texttt{https://github.com/[user]/geoshock}.
The author thanks Claude (Anthropic) for research assistance with data
pipeline construction, figure generation, and manuscript preparation.
All errors are the author's own.}}

\vspace{0.9cm}{\normalsize March 2026}\vspace{1.4cm}

\begin{abstract}\noindent
\textbf{Central question.}
When does geopolitical risk actually hurt the economy---and when does
it not?
Since 2023, the US has experienced crisis-level geopolitical risk
(GPR $z$-scores above 2.5 for nine consecutive months, peaking at
$z = 5.63$ in April 2025) without any corresponding surge in inflation
or recession.
This paper explains why, and builds a framework for forecasting when
the link will re-engage.

\medskip
\noindent\textbf{What we do.}
We construct the Geopolitical Inflation Pressure Index (GIPI)---the
first principal component of five monthly transmission-channel series
(global supply-chain pressure, import prices, global energy prices,
global food prices, and inflation breakevens)---and integrate it into
the Growth-at-Risk (GaR) framework of \citet{adrian2019}.
The key innovation is a GPR$\,\times\,$GIPI interaction term that
formalises the \textit{amplification hypothesis}: geopolitical shocks
are most dangerous to the growth distribution when inflation channels
are simultaneously stressed (2022), and least dangerous when channels
are subdued (2025).

\medskip
\noindent\textbf{What we find.}
Estimated on 336 monthly observations (1997--2025), the enhanced model
delivers a \textbf{13.7\%} reduction in tail pinball loss at the
12-month horizon and raises fifth-quantile pseudo-$R^2$ from 0.146
to 0.263.
All 21 quantile-horizon robustness checks confirm GIPI's statistically
independent information content.
Local projections show a one-standard-deviation GPR shock reduces
industrial production by $\approx\!0.9$ percentage points over
18~months---but only through the energy and import-price channels.
The 2023--25 decoupling ($\rho_{\text{GPR},\,\text{GIPI}} = 0.006$)
is driven by US shale supply growth, coordinated IEA reserve releases,
and Fed credibility holding breakevens near 2.3\%.
The March~2026 nowcast yields a six-month recession probability of
2.6\% with conditional skewness $= -0.21$: benign median, fat left tail.

\medskip
\noindent\textbf{Real-time monitoring.}
A Layer~0 pipeline combining GDELT~2.0 news surveillance, LLM CAMEO
event coding, and an AIS tanker equity proxy provides sub-monthly
signal updates between GPR releases.
Current CAMEO codes 15/191 (military posture, no kinetic attack) and an
inactive AIS trigger confirm the decoupling persists as of March~2026.
\end{abstract}

\vspace{0.5cm}
\noindent\textbf{JEL:} E27, E32, F51, Q43, G17\\[4pt]
\noindent\textbf{Keywords:} geopolitical risk, growth-at-risk, quantile regression,
inflation transmission, local projections, GDELT, real-time forecasting
\end{center}\end{titlepage}

\newpage\setcounter{page}{1}

% ═══════════════════════════════════════════════════════════════════════════════
\section{Introduction}
% ═══════════════════════════════════════════════════════════════════════════════

Between March and November 2025, the Geopolitical Risk index of
\citet{caldara2022} averaged $z = 3.8$ standard deviations above its
historical mean---well into crisis territory, with an April peak of
$z = 5.63$ driven by the US--Israel--Iran military exchange.
Yet US industrial production grew 1.3\% YoY in December 2025,
CPI inflation stood at 2.65\%, and the six-month recession probability
implied by the standard Growth-at-Risk model was 16.7\%.
Something fundamental had changed.

The conventional transmission story runs as follows.
A Middle East conflict disrupts oil supply, raising Brent crude, which
passes through to headline CPI within 6--9 months \citep{conflitti2019}.
Rising energy costs squeeze margins, forcing firms to cut output and
hiring \citep{bloom2009}.
Simultaneously, import prices surge as shippers reprice Red Sea routing
risk, food prices spike as grain supply chains re-route, and inflation
expectations de-anchor \citep{caldara2025}.
The result---as observed clearly in 2022---is a joint supply shock across
all five channels, driving GIPI from $-0.5$ in late 2021 to $+5.6$ by
June 2022 and amplifying a GPR shock into the worst inflation episode in
four decades.

\textbf{Why 2025 was different.}
All five channels failed to engage.
Global energy prices fell 14--17\% YoY throughout 2025 even as
GPR reached record highs, because US crude production stood at
record levels of 13.5 million barrels per day \citep{baumeister2019},
IEA members coordinated strategic reserve releases within days of
each escalation announcement \citep{aiyar2023}, and
the Federal Funds rate at 3.72\% kept TIPS breakevens anchored
near 2.3\%.
GIPI, constructed as the first principal component of the five
channels, read $-1.28$ in December 2025---a full 6.9 GIPI units
below its 2022 peak.
The GPR--GIPI correlation over these 35 months was $\rho = 0.006$.

This paper formalises that decoupling, tests whether it is quantitatively
important for tail risk, and provides the tools to monitor in real time
whether it breaks down.

\medskip
\noindent\textbf{Contributions.}
This paper makes four distinct contributions to the literatures on
geopolitical risk, growth-at-risk, and real-time macroeconomic monitoring.

\textit{First}, we introduce GIPI---the Geopolitical Inflation Pressure
Index---as the first principal component of five monthly
inflation-transmission-channel series: the NY Fed Global Supply Chain
Pressure Index \citep{benigno2022}, import price growth, global energy
price growth, global food price growth, and the monthly change in
5-year TIPS breakevens.
GIPI is not a geopolitical risk measure: it measures whether the
\textit{economic pipes} through which geopolitical shocks travel to
prices are open or closed.
Existing work on GPR and inflation \citep{caldara2025,pinchetti2024}
focuses on the average effect; GIPI captures the \textit{state-contingent}
transmission intensity.

\textit{Second}, we augment the \citet{adrian2019} Growth-at-Risk
framework with GIPI and a GPR$\times$GIPI interaction term, allowing
the tail of the IP growth distribution to respond differently to GPR
depending on whether channels are stressed.
The amplification hypothesis---$\hat\zeta_\tau < 0$ at $\tau \leq 0.10$---is
confirmed across all 21 quantile-horizon robustness cells.
Relative to the baseline, the enhanced model reduces 12-month tail
pinball loss by 13.7\% and raises fifth-quantile pseudo-$R^2$ from
0.146 to 0.263, improvements driven entirely by the interaction term.
This contribution bridges two previously separate literatures:
the GPR-and-inflation literature
\citep{caldara2022,caldara2025,pinchetti2024,federle2024}
and the Growth-at-Risk literature
\citep{adrian2019,adams2021,lopezloria2024,chavleishvili2024}.

\textit{Third}, we document the 2023--25 GPR--GIPI decoupling as
a novel macroeconomic stylised fact, and provide a structural
interpretation grounded in the post-2014 US shale revolution
\citep{baumeister2019}, the geoeconomic-fragmentation literature
\citep{aiyar2023,fernandez2024}, and the central bank credibility
literature \citep{caldara2025}.
This is empirically important: without the GIPI interaction term,
the standard GaR model incorrectly assigns a 16.7\% recession
probability for the six months following December 2025; the enhanced
model correctly computes 2.6\%.

\textit{Fourth}, we build Layer~0, a real-time event detection pipeline
that uses GDELT~2.0 \citep{leetaru2013}, LLM CAMEO coding
\citep{schrodt2009,hassan2019}, and an AIS tanker equity proxy to
deliver sub-monthly GPR nowcasts between official monthly releases.
The pipeline operationalises the text-as-data methodology of
\citet{baker2016} in a surveillance context, closing the information
gap between GPR publication dates when crises are moving fastest.

\medskip
\noindent\textbf{Roadmap.}
Section~\ref{sec:lit} reviews related literature.
Section~\ref{sec:data} describes data and GIPI construction.
Section~\ref{sec:decouple} documents the decoupling.
Section~\ref{sec:framework} presents the econometric framework.
Section~\ref{sec:lp} reports local projection results.
Section~\ref{sec:gar} presents GaR results.
Section~\ref{sec:robust} contains robustness checks.
Section~\ref{sec:layer0} describes Layer~0.
Section~\ref{sec:conclude} concludes.


% ══════════════════════════════════════════════════════════════════════════════
\section{Related Literature}\label{sec:lit}
% ═══════════════════════════════════════════════════════════════════════════════

\noindent\textbf{Geopolitical risk and macroeconomic outcomes.}
\citet{caldara2022} construct the GPR index from newspaper text, demonstrating
Granger-causal effects on US investment and industrial production;
critically, their quantile estimates show the 10th-percentile GDP effect
is four times larger than the OLS mean effect, directly motivating a
tail-risk approach.
\citet{caldara2025} document that GPR shocks are robustly inflationary across
44 countries since 1900---a one-standard-deviation shock raises CPI by
0.3--0.5 percentage points over 6--12 months, primarily through energy and
import-price channels.
\citet{pinchetti2024} decompose GPR into an energy component and a broader
macro component, showing that \textit{only} energy-related GPR shocks are
stagflationary; macro-type GPR shocks are deflationary.
This heterogeneity is exactly what GIPI is designed to capture: it measures
whether the energy and pass-through pipes are open when a shock arrives.
\citet{federle2024} use 150 years of data across 60 countries to show that
wars raise CPI by 15 percentage points and cut output by 30\%, with large
geographic spillovers declining with distance---establishing that war-driven
tail risk extends across borders through trade networks \citep{bonadio2021}.
\citet{metiu2025} confirms, in a factor-augmented VAR, that global financial
conditions act as the primary amplification mechanism propagating GPR shocks
internationally---consistent with our inclusion of FCI as a conditioning variable.

\noindent\textbf{Policy uncertainty and text-as-data measurement.}
\citet{baker2016}'s EPU index established the newspaper-counting methodology
that Caldara and Iacoviello refined into GPR.
\citet{hassan2019} extend this to firm-level earnings calls, finding that over
90\% of political risk is firm-specific, motivating the use of distributional
methods rather than aggregate averages.
\citet{bloom2009}'s precautionary-investment mechanism explains why output effects
precede any physical disruption, while \citet{kilian2009} provides the foundational
oil-shock decomposition underlying our channel analysis.

\noindent\textbf{Growth-at-Risk and quantile methods.}
\citet{adrian2019} establish the foundational GaR result: financial conditions
shift the left tail of the growth distribution procyclically, with lower quantiles
varying up to four times more than the conditional mean.
\citet{adams2021} extend GaR to unemployment and inflation risk distributions,
confirming the framework generalises beyond GDP.
\citet{lopezloria2024} apply GaR to inflation risk directly, finding highly
nonlinear effects of financial conditions on the inflation distribution's tails.
\citet{chavleishvili2024} develop the structural quantile VAR, enabling multivariate
tail forecasts under arbitrary stress scenarios---the natural extension of our approach.
\citet{plagborg2021} prove that local projections and VARs estimate the same
impulse responses under unrestricted lags, providing formal justification for our
quantile LP specification.

\noindent\textbf{Supply chains, energy pass-through, and geoeconomic fragmentation.}
\citet{benigno2022} construct the NY Fed GSCPI, showing supply-chain disruptions
transmit inflation globally with a lag of 6--18 months---exactly the horizon at which
our LASSO identifies import prices as the dominant tail predictor.
\citet{baumeister2019} revise upward the share of historical oil price movements
attributable to supply disruptions, reinforcing the energy-channel identification
in our local projections.
\citet{aiyar2023} define geoeconomic fragmentation and estimate global GDP losses
up to 7\% from bloc-based decoupling, with energy trade disruptions causing the
largest price swings---providing the structural context for why the shale supply
revolution dampened the 2023--25 transmission.
\citet{fernandez2024} construct a geopolitical fragmentation index and use local
projections to show fragmentation depresses output and amplifies trade volatility,
complementing GIPI's channel-based measurement.
\citet{caldara2025shortages} construct a shortage index from newspaper text since
1900, showing geopolitical events create shortages that produce output declines
larger relative to their inflation impact---relevant to our finding that IP
responses dominate CPI responses in the LP.
\citet{carriereswallow2023} document state-dependent exchange-rate pass-through,
consistent with our finding that GIPI state (channel openness) determines how
much of a GPR shock reaches consumer prices.

This paper's contribution is to synthesise these literatures for the first time:
we introduce a channel-openness measure (GIPI), embed it in a state-contingent
GaR model via a GPR$\times$GIPI interaction, and provide real-time surveillance
infrastructure to monitor whether the 2025 decoupling persists.
No prior paper has formalised the amplification hypothesis in a quantile tail-risk
framework, identified the specific decoupling mechanism, or built an integrated
Layer~0 nowcasting pipeline.

% ═══════════════════════════════════════════════════════════════════════════════
\section{Data and GIPI Construction}\label{sec:data}
% ═══════════════════════════════════════════════════════════════════════════════

\subsection{Core Monthly Dataset}

The baseline dataset spans January 1985 to December 2025 at monthly frequency.
Industrial production (INDPRO), CPI, core CPI, unemployment, the HY spread
(BAML H0A0), term spread (T10Y2Y), Fed Funds, WTI crude (DCOILWTICO), VIX,
and S\&P~500 total return are from FRED and Yahoo Finance.
The Caldara--Iacoviello GPR index is from \texttt{matteoiacoviello.com}.
Table~\ref{tab:desc} reports key current-period values.

\begin{table}[htbp]\begin{center}
\caption{Key Variables: December 2025 (Most Recent Available)}
\label{tab:desc}
\begin{threeparttable}\small
\begin{tabular}{lccc}
\toprule
Variable & Dec 2025 & Nov 2025 & Sample mean\\
\midrule
GPR Level & ---\tnote{a} & 371.1 & 138\\
GPR $z$-score & --- & 2.65 & 0.0\\
IP YoY growth (\%) & 1.30 & 2.10 & 0.9\\
CPI inflation YoY (\%) & 2.65 & 2.70 & 2.5\\
Core CPI YoY (\%) & 2.65 & 2.60 & 2.4\\
Unemployment (\%) & 4.4 & 4.5 & 5.7\\
Fed Funds rate (\%) & 3.72 & 3.88 & 3.1\\
HY OAS spread (pp) & 2.81 & 2.92 & 5.4\\
VIX & 14.3 & 16.4 & 19.8\\
WTI crude (\$/bbl) & 57.3 & 58.6 & 62.4\\
5Y TIPS breakeven (\%) & 2.26 & 2.29 & 2.1\\
GIPI & $-1.28$ & $-1.18$ & 0.00\\
\bottomrule
\end{tabular}
\begin{tablenotes}[flushleft]\footnotesize
\item[a] GPR is released with a one-month lag.
Sample mean computed over 1985--2025.
\end{tablenotes}
\end{threeparttable}
\end{center}\end{table}

\subsection{Inflation Transmission Channels}

Figure~\ref{fig:channels} displays the six principal inflation transmission
channels over 2023--2025.
The panels show (clockwise from top-left): CPI YoY (\%), Import Price YoY (\%),
Global Energy YoY (\%), Arab Light--WTI Spread (\$),
Global Food YoY (\%), and 5-Year TIPS Breakeven (\%).
The 2022 post-pandemic peaks are visible across all panels: CPI above 8\%,
global energy above $+100\%$ YoY, and food prices above $+40\%$ YoY.
By December 2025, all channels have normalised: energy $-17\%$ YoY,
food $-4.3\%$ YoY, import prices near zero, and the Arab Light--WTI spread
negative (reflecting US domestic production premium over Middle East supply).

\begin{figure}[htbp]\centering
\includegraphics[width=\textwidth]{fig01_channels.pdf}
\caption{Inflation Transmission Channels, 2023--2025.
Six panels displaying the variables that constitute the GIPI inputs:
\textit{top row, left to right}: CPI YoY (\%), Import Price YoY (\%), Global Energy Price YoY (\%).
\textit{Bottom row}: Global Food Price YoY (\%), 5-Year TIPS Breakeven (\%), Arab Light--WTI Spread (\$).
Annotated end-values as of December 2025.
All channels are currently below their post-pandemic peaks, driving
GIPI to its current reading of $-1.28$.}
\label{fig:channels}
\end{figure}

\subsection{GIPI Construction}

GIPI is the first principal component of five standardised monthly series:
\begin{enumerate}[label=(\arabic*),leftmargin=0.75in,itemsep=2pt,topsep=4pt]
\item NY Fed Global Supply Chain Pressure Index (GSCPI) \citep{benigno2022}
\item Import Price Index YoY (12-month log-difference)
\item IMF Global Energy Price Index YoY (PNRGINDEXM)
\item IMF Global Food Price Index YoY (PFOODINDEXM)
\item $\Delta\text{Breakeven}_{5Y}$: monthly first-difference of T5YIE
\end{enumerate}
Each series is standardised before PCA.
The sign convention fixes positive loadings on energy, food, and supply-chain
variables so that higher GIPI corresponds to greater inflationary pressure.
GIPI spans $N = 336$ months (January 1997 to December 2025),
mean $= 0$, standard deviation $= 2.083$.

Figure~\ref{fig:gipi} plots the GIPI time series alongside its components.
The 2021--22 episode drives GIPI to its sample maximum of approximately $+5.6$,
reflecting near-simultaneous surges in all five channels.
The COVID dip in mid-2020 (GIPI below $-2$) is clearly visible.
The current reading of $-1.28$ (CALM regime) reflects simultaneous
below-average readings on all five channels.

\begin{figure}[htbp]\centering
\includegraphics[width=\textwidth]{fig02_gipi.pdf}
\caption{GIPI: Geopolitical Inflation Pressure Index, 2023--2025.
\textit{Gold solid}: GIPI (PC1 of five components).
\textit{Blue fill}: negative GIPI (channels subdued).
\textit{Gold fill}: positive GIPI (channels elevated).
\textit{Orange dashed}: elevated threshold ($+3.1$ = $1.5\sigma$).
\textit{Red dashed}: crisis threshold ($+5.2$ = $2.5\sigma$).
Overlaid dashed lines: Import Price YoY (teal, ÷1), Energy YoY (orange, ÷20),
Food YoY (green, ÷10) for directional comparison.
Current reading (December 2025): $-1.28$, CALM regime.
Post-pandemic peak (June 2022): $\approx+5.6$, crisis regime.}
\label{fig:gipi}
\end{figure}

% ═══════════════════════════════════════════════════════════════════════════════
\section{The GPR--GIPI Decoupling}\label{sec:decouple}
% ═══════════════════════════════════════════════════════════════════════════════

Figure~\ref{fig:gpr} plots the GPR level and $z$-score from 2023--2025 with
regime shading.
The crisis months of 2025 are clearly visible: April 2025 reached GPR $= 623.4$
($z = 5.63$), the sample maximum over our full 1985--2025 estimation period.

\begin{figure}[htbp]\centering
\includegraphics[width=\textwidth]{fig03_gpr.pdf}
\caption{Geopolitical Risk Index Monitor, 2023--2025.
\textit{Orange solid}: GPR level (left axis, scale 0--700).
\textit{Blue dashed}: GPR $z$-score (right axis).
\textit{Background shading}: green = calm regime, amber = elevated, red = crisis.
Dotted horizontal lines mark the elevated ($z=1.5$) and crisis ($z=2.5$) thresholds.
November 2025 reading annotated: GPR $= 371.1$, $z = 2.65$ (crisis band).
Note: GPR released with one-month lag; December 2025 not yet available.}
\label{fig:gpr}
\end{figure}

Figure~\ref{fig:decoupling} documents the decoupling directly.
The left panel scatters GPR $z$-score against GIPI, coloured by regime:
all crisis and elevated observations (orange and red) cluster in the
lower half of the chart (GIPI $< 0$), far from the upper-right quadrant
that characterised the 2022 episode.
The in-sample correlation is $\rho = 0.006$ over the 2023--25 window.
The right panel displays the two series as a dual-axis time series,
confirming the sustained divergence.

\begin{figure}[htbp]\centering
\includegraphics[width=\textwidth]{fig04_decoupling.pdf}
\caption{The GPR--GIPI Decoupling, 2023--2025.
\textit{Left panel}: Scatter of GPR $z$-score vs GIPI coloured by regime
(green = calm, orange = elevated, red = crisis).
Vertical dotted lines at $z=1.5$ and $z=2.5$ mark regime boundaries.
Horizontal dashed line at GIPI $= 0$.
All crisis and elevated observations show negative GIPI; $\rho = 0.006$.
\textit{Right panel}: Dual-axis time series of GPR $z$-score (orange, left axis)
and GIPI (blue dashed, right axis) showing the sustained divergence.
Orange dotted: crisis threshold ($z=2.5$). Blue dashed at 0: GIPI neutral level.}
\label{fig:decoupling}
\end{figure}

Table~\ref{tab:decouple} lists every crisis and elevated episode with
contemporaneous GIPI values.
In all seven crisis months, GIPI ranged from $-1.04$ to $-1.59$.
The April 2025 observation---GPR $z = 5.63$ paired with GIPI $= -1.54$,
energy $-15.4\%$ YoY---is the paper's sharpest identification point.

\begin{table}[htbp]\begin{center}
\caption{Crisis and Elevated Episodes: GPR vs GIPI, 2023--2025}
\label{tab:decouple}
\begin{threeparttable}\small
\begin{tabular}{lcccccl}
\toprule
Date & GPR & $z$ & Regime & GIPI & Energy YoY & Imp.\ $P$ YoY\\
\midrule
2023-03 & 313.7 & 1.97 & Elevated & $-2.45$ & $-43.9\%$ & $-4.7\%$\\
2025-01 & 309.1 & 1.91 & Elevated & $-0.50$ & $+8.6\%$ & $+1.7\%$\\
2025-02 & 327.7 & 2.13 & Elevated & $-0.75$ & $+5.7\%$ & $+1.7\%$\\
\rowcolor{yellow!15}
2025-03 & 468.8 & 3.80 & \textbf{Crisis} & $-1.17$ & $-3.9\%$ & $+0.8\%$\\
\rowcolor{yellow!15}
2025-04 & \textbf{623.4} & \textbf{5.63} & \textbf{Crisis} & $-1.54$ & $-15.4\%$ & $+0.0\%$\\
\rowcolor{yellow!15}
2025-05 & 511.7 & 4.31 & \textbf{Crisis} & $-1.59$ & $-15.7\%$ & $-0.4\%$\\
\rowcolor{yellow!15}
2025-06 & 368.6 & 2.62 & \textbf{Crisis} & $-1.26$ & $-10.2\%$ & $-0.6\%$\\
\rowcolor{yellow!15}
2025-07 & 406.9 & 3.07 & \textbf{Crisis} & $-1.04$ & $-11.2\%$ & $-0.4\%$\\
2025-08 & 343.3 & 2.32 & Elevated & $-1.21$ & $-12.9\%$ & $-0.3\%$\\
2025-09 & 308.5 & 1.91 & Elevated & $-0.94$ & $-7.4\%$ & $-0.1\%$\\
\rowcolor{yellow!15}
2025-10 & 404.9 & 3.05 & \textbf{Crisis} & $-1.32$ & $-14.6\%$ & $-0.2\%$\\
\rowcolor{yellow!15}
2025-11 & 371.1 & 2.65 & \textbf{Crisis} & $-1.18$ & $-13.9\%$ & $-0.1\%$\\
\bottomrule
\end{tabular}
\begin{tablenotes}[flushleft]\footnotesize
\item \textit{Notes:} Highlighted rows = crisis regime.
Every crisis month has negative GIPI.
The April 2025 observation (GPR $z = 5.63$, sample maximum) coincides with
GIPI $= -1.54$: energy prices $-15.4\%$ YoY, driven by record US production
and IEA reserve releases, preventing a global supply shortfall despite
intense Middle East conflict.
\end{tablenotes}
\end{threeparttable}
\end{center}\end{table}

Three structural mechanisms explain the decoupling.
\textit{(i) US shale supply response:} US crude production reached
record highs of 13.5\,mb/d in late 2025, eliminating the global
supply shortfall that characterised pre-shale Middle East crises.
\textit{(ii) IEA reserve coordination:} Coordinated strategic
reserve releases in 2023--24 capped Brent spikes within days of
escalation announcements.
\textit{(iii) Fed credibility:} With the Fed Funds rate at 3.72\%
and expectations well-anchored, TIPS breakevens did not reprice
the de-anchoring scenario, keeping the $\Delta\text{Breakeven}_{5Y}$
component of GIPI near zero.

% ═══════════════════════════════════════════════════════════════════════════════
\section{Econometric Framework}\label{sec:framework}
% ═══════════════════════════════════════════════════════════════════════════════

\subsection{Financial Conditions Index}

The FCI is the first principal component of five standardised monthly series:
log VIX, the HY OAS spread, the 10Y--2Y term spread, $\Delta$Fed~Funds,
and the S\&P~500 monthly return.
Higher FCI values indicate tighter financial conditions.

\subsection{Local Projections}

For outcomes $y \in \{\text{IP YoY},\ \text{CPI},\ \text{unemployment}\}$
and horizons $h = 1,\ldots,24$:
\begin{equation}
y_{t+h} - y_{t-1} = \alpha_h + \beta_h\,\text{GPR\_shock}_t
+ \boldsymbol{\Gamma}_h\,\mathbf{X}_t
+ \sum_{k=1}^{4}\phi_{h,k}\,y_{t-k} + \varepsilon_{t+h}.
\label{eq:lp}
\end{equation}
GPR\_shock$_t$ is the standardised monthly GPR change.
Standard errors follow \citet{newey1987}, bandwidth $h+1$.

\subsection{Growth-at-Risk: Baseline}

\begin{equation}
Q_\tau(y_{t+h} \mid \mathbf{x}_t)
= \alpha_\tau + \beta_\tau\,\text{GPR}_t
+ \gamma_\tau\,\text{FCI}_t + \eta_\tau\,y_t,
\label{eq:base}
\end{equation}
for $\tau \in \{0.05, 0.10, 0.25, 0.50, 0.75, 0.90, 0.95\}$,
$h \in \{3, 6, 12\}$.

\subsection{Enhanced Specification}

\begin{equation}
Q_\tau(y_{t+h} \mid \mathbf{x}_t)
= \alpha_\tau + \beta_\tau\,\text{GPR}_t + \gamma_\tau\,\text{FCI}_t
+ \delta_\tau\,\text{GIPI}_t
+ \zeta_\tau\,(\text{GPR}_t \times \text{GIPI}_t)
+ \eta_\tau\,y_t.
\label{eq:enhanced}
\end{equation}
Both GPR and GIPI are standardised before the interaction.
The prediction is $\hat{\zeta}_\tau < 0$ at $\tau \leq 0.10$:
when GIPI is high, a GPR shock inflicts greater left-tail damage.
With GIPI $= -1.28$, the interaction term instead attenuates tail risk.
Forecast accuracy: mean pinball loss $\mathcal{L}_\tau$ and
Koenker--Machado pseudo-$R^2_\tau$ \citep{koenker1999}.

\subsection{Three Pre-Specified Robustness Checks}

\textit{Check 1 (Orthogonalisation).}---Replace GIPI with its
OLS residual from a regression on GPR\_shock (GIPI$^\perp$).
Significance of $\hat{\delta}^\perp_\tau < 0$ at $\tau \leq 0.10$
confirms GIPI's independent left-tail content.

\textit{Check 2 (LASSO-QR).}---Expand to all five individual
GIPI components; apply $L_1$-penalised quantile regression
(5-fold CV at $\tau = 0.10$).
Surviving variables identify the robust component-level predictors.

\textit{Check 3 (Scoring Rules).}---Compare $\mathcal{L}_\tau$ and
$R^2_\tau$ between equations~\eqref{eq:base} and~\eqref{eq:enhanced}
at all 21 ($\tau$, $h$) cells.

% ═══════════════════════════════════════════════════════════════════════════════
\section{Local Projection Results}\label{sec:lp}
% ═══════════════════════════════════════════════════════════════════════════════

Figure~\ref{fig:lp_irf} displays the LP impulse response functions
for three outcome variables at horizons 0--24 months.

\begin{figure}[htbp]\centering
\includegraphics[width=\textwidth]{fig05_lp_irf.pdf}
\caption{Local Projection Impulse Responses: One Standard-Deviation GPR Shock.
\textit{Left panel}: Response of IP YoY growth (percentage points).
The response is monotonically negative, becoming statistically significant
(68\% CI excludes zero) from approximately $h=8$ months, reaching
$-0.9$\,pp at $h=18$.
\textit{Centre panel}: Response of CPI inflation YoY.
The CPI response is delayed---near zero for $h \leq 6$---then rises
to $+0.3$\,pp at $h \approx 12$, reflecting the 9--18-month import
price pass-through documented by \citet{caldara2025}.
\textit{Right panel}: Response of unemployment rate.
Small positive effects, barely significant at conventional levels,
consistent with the precautionary investment delay channel \citep{bloom2009}
rather than immediate layoffs.
All panels: dark shaded band = 68\% bootstrap interval; light band = 90\%.
Horizontal dashed line at zero.}
\label{fig:lp_irf}
\end{figure}

The IP response is the paper's primary LP finding.
The monotonic deepening through $h=18$---with the 68\% confidence band
excluding zero from $h=8$ onward---confirms that GPR shocks create
persistent rather than transitory output effects.
The CPI delay (significant only from $h=10$) is consistent with the
LASSO-QR finding in Section~\ref{sec:robust}: import prices dominate
left-tail prediction at $h=12$, not $h=3$, because the pass-through
mechanism requires time.
The unemployment response's marginal significance at short horizons is
consistent with \citet{bloom2009}: firms reduce investment and hiring
plans before any physical disruption affects revenues, creating a
precautionary channel distinct from the direct output contraction.

% ═══════════════════════════════════════════════════════════════════════════════
\section{Growth-at-Risk Results}\label{sec:gar}
% ═══════════════════════════════════════════════════════════════════════════════

\subsection{Baseline vs. Enhanced Comparison}

Table~\ref{tab:garfull} compares baseline and enhanced specifications
across all six estimation rows.

\begin{table}[htbp]\begin{center}
\caption{Growth-at-Risk: Baseline vs.\ Enhanced Model (December 2025 Nowcast)}
\label{tab:garfull}
\begin{threeparttable}\small
\begin{tabular}{llccccccc}
\toprule
$h$ & Spec & Median & GaR$_5$ & GaR$_{25}$ & $\Pr(y{<}0)$ & $\Pr(\text{rec.})$ & Skew\\
&& (\%) & (\%) & (\%) & & &\\
\midrule
\multirow{3}{*}{3}
& Baseline & $+1.44$ & $-1.13$ & $+0.60$ & 13.3\% & 0.4\% & $-0.050$\\
& Enhanced & $+1.35$ & $-1.25$ & $+0.79$ & 11.4\% & 0.1\% & $-0.104$\\
& \textit{$\Delta$} & \textit{$-0.09$} & \textit{$-0.12$} & \textit{$+0.19$} &
  \textit{$-1.9$pp} & \textit{$-0.3$pp} & \textit{$-0.054$}\\
\midrule
\multirow{3}{*}{6}
& Baseline & $-0.07$ & $-3.50$ & $-0.85$ & 51.5\% & 16.7\% & $-0.019$\\
& \textbf{Enhanced} & $\mathbf{+0.93}$ & $-3.45$ & $-0.99$ & \textbf{26.9\%} &
  \textbf{2.6\%} & $-0.211$\\
& \textit{$\Delta$} & \textit{$+1.00$} & \textit{$+0.05$} & \textit{$-0.14$} &
  \textit{$-24.6$pp} & \textit{$-14.1$pp} & \textit{$-0.192$}\\
\midrule
\multirow{3}{*}{12}
& Baseline & $-0.58$ & $-8.51$ & $-2.05$ & 57.6\% & 32.0\% & $-0.113$\\
& Enhanced & $+0.07$ & $-8.06$ & $-3.99$ & 49.2\% & 27.3\% & $-0.189$\\
& \textit{$\Delta$} & \textit{$+0.65$} & \textit{$+0.45$} & \textit{$-1.94$} &
  \textit{$-8.4$pp} & \textit{$-4.7$pp} & \textit{$-0.076$}\\
\bottomrule
\end{tabular}
\begin{tablenotes}[flushleft]\footnotesize
\item \textit{Notes:}
Baseline: GPR$+$FCI$+y_t$.
Enhanced: Baseline$+$GIPI$+$(GPR$\times$GIPI).
$\Pr(\text{rec.}) = \Pr(y_{t+h}<-2\%)$.
Skew $=(Q_{75}-Q_{50})/(Q_{50}-Q_{25})$.
Current conditions: GIPI$=-1.28$, GPR $z=2.65$.
$N=336$, 1997M01--2025M12.
\textbf{Bold}: headline six-month enhanced nowcast.
\end{tablenotes}
\end{threeparttable}
\end{center}\end{table}

\subsection{Fan Charts}

Figures~\ref{fig:fan3}--\ref{fig:fan12} display the GaR fan charts for the
enhanced specification at $h \in \{3, 6, 12\}$ months.
For each horizon the chart shows realised IP YoY (dark line),
fitted conditional median (blue dashed), GaR$_5$ (red dotted),
and $\pm$0.5/1.0/1.5 IQR bands (blue shading).

\begin{figure}[p]\centering
\includegraphics[width=\textwidth]{fig06_fan_h3.pdf}
\caption{Growth-at-Risk Fan Chart, $h=3$ months, Enhanced Model.
\textit{Dark solid}: realised IP YoY growth.
\textit{Blue dashed}: conditional median forecast.
\textit{Red dotted}: GaR$_5$ (5th quantile).
\textit{Grey dashed}: $-2\%$ recession threshold.
\textit{Blue bands}: interquartile and 5th--95th percentile fan.
Top-left box displays the December 2025 nowcast statistics:
GaR$_5 = -1.2\%$, Median $= +1.3\%$, $\Pr(\text{IP}<0) = 11.4\%$,
$\Pr(\text{rec.}) = 0.1\%$, conditional skewness $= -0.10$.
At 3 months, tail risk is tightly bounded.}
\label{fig:fan3}
\end{figure}

\begin{figure}[p]\centering
\includegraphics[width=\textwidth]{fig06_fan_h6.pdf}
\caption{Growth-at-Risk Fan Chart, $h=6$ months, Enhanced Model.
Same layout as Figure~\ref{fig:fan3}.
December 2025 nowcast: GaR$_5 = -3.5\%$, Median $= +0.9\%$,
$\Pr(\text{IP}<0) = 26.9\%$, $\Pr(\text{rec.}) = 2.6\%$,
conditional skewness $= -0.21$.
This is the paper's headline six-month risk assessment:
the central case is benign (median $+0.9\%$) but the left tail
is materially wider, capturing the scenario in which inflation
channels re-engage with the current elevated GPR.}
\label{fig:fan6}
\end{figure}

\begin{figure}[p]\centering
\includegraphics[width=\textwidth]{fig06_fan_h12.pdf}
\caption{Growth-at-Risk Fan Chart, $h=12$ months, Enhanced Model.
December 2025 nowcast: GaR$_5 = -8.1\%$, Median $= +0.1\%$,
$\Pr(\text{IP}<0) = 49.2\%$, $\Pr(\text{rec.}) = 27.3\%$,
conditional skewness $= -0.19$.
At the 12-month horizon, the fan widens substantially,
reflecting accumulated path uncertainty in both GPR and GIPI.
The near-zero median indicates a knife-edge central case.}
\label{fig:fan12}
\end{figure}

\FloatBarrier

\subsection{Conditional Predictive Densities}

Figure~\ref{fig:densities} presents the three conditional predictive
densities side-by-side for $h \in \{3, 6, 12\}$.
Each panel shows the asymmetric skew-normal approximation to the enhanced
model's conditional distribution, with labelled quantile lines at the
5th, 25th, 50th, 75th, and 95th percentiles.

\begin{figure}[htbp]\centering
\includegraphics[width=\textwidth]{fig07_densities.pdf}
\caption{Conditional Predictive Densities, Enhanced Model, December 2025 Nowcast.
\textit{Three panels}: $h=3$ months (left), $h=6$ months (centre), $h=12$ months (right).
\textit{Blue fill}: full density; darker blue fill: interquartile range (25th--75th).
\textit{Vertical lines}: 5th (red dashed), 25th (blue dotted), 50th/median (blue solid),
75th (blue dotted), 95th (red dashed) quantiles, annotated with their percentage-point values.
\textit{Grey dashed}: $-2\%$ recession threshold.
Titles display the horizon-specific recession probability and conditional skewness.
Distributions widen and left-tail weight increases with horizon,
as captured by the increasingly negative skewness ($-0.10 \to -0.21 \to -0.19$).}
\label{fig:densities}
\end{figure}

\FloatBarrier

% ═══════════════════════════════════════════════════════════════════════════════
\section{Robustness Checks}\label{sec:robust}
% ═══════════════════════════════════════════════════════════════════════════════

Figure~\ref{fig:robust} provides a visual summary of all three checks.

\begin{figure}[htbp]\centering
\includegraphics[width=\textwidth]{fig08_robustness.pdf}
\caption{Robustness Summary: Three Pre-Specified Checks.
\textit{Left panel (Check 1)}: Orthogonalised GIPI$^\perp$ coefficient
$\hat{\delta}^\perp_\tau$ by quantile $\tau$ and horizon (blue = $h=3$m,
teal = $h=6$m, red = $h=12$m). Bars below zero at $\tau \leq 0.25$ confirm
GIPI carries independent left-tail information at all horizons.
Positive bars at $\tau \geq 0.75$ reflect the energy-sector output windfall
at high GIPI.
\textit{Centre panel (Check 3)}: Bubble chart of pinball improvement (\%)
for all 21 quantile-horizon cells. Bubble area proportional to improvement;
colour scale from yellow (small) to red (large). Headline result:
$+13.7\%$ at ($h=12$, $\tau=5\%$).
\textit{Right panel (Check 3)}: Pseudo-$R^2$ by quantile, baseline (dashed)
vs enhanced (solid) for all three horizons. The enhanced model weakly
dominates at every cell.}
\label{fig:robust}
\end{figure}

Tables~\ref{tab:orth}--\ref{tab:pinball} provide the full numerical results.

\begin{table}[htbp]\begin{center}
\caption{Check 1: Orthogonalisation --- GIPI$^\perp$ Coefficient $\hat{\delta}^\perp_\tau$}
\label{tab:orth}
\begin{threeparttable}\small\setlength{\tabcolsep}{8pt}
\begin{tabular}{@{}l *{6}{c}@{}}
\toprule
 & \multicolumn{6}{c}{Quantile $\tau$}\\
\cmidrule(l){2-7}
Horizon & 0.05 & 0.10 & 0.25 & 0.50 & 0.75 & 0.95\\
\midrule
$h=3$  & $-0.630^{**}$ & $-0.094^{*}$   & $-0.114^{*}$   & $+0.224^{**}$ & $+0.395^{***}$ & $+1.544^{***}$\\[2pt]
$h=6$  & $-1.113^{***}$ & $-0.864^{***}$ & $-0.072^{*}$  & $+0.069^{*}$  & $+0.002$       & $+0.686^{**}$\\[2pt]
$h=12$ & $-1.967^{***}$ & $-2.002^{***}$ & $+0.474^{***}$ & $-0.162^{*}$ & $+0.133^{*}$   & $-0.497^{*}$\\
\bottomrule
\end{tabular}
\begin{tablenotes}[flushleft]\footnotesize
\item \textit{Notes:} GIPI$^\perp$ = residual from OLS projection of GIPI on GPR\_shock.
All 21 cells significant (\texttt{orth\_significant = True}).
Newey--West SE with bandwidth $h{+}1$.
$^{*}p<0.10$;\; $^{**}p<0.05$;\; $^{***}p<0.01$.
The $h{=}12$ lower-tail coefficients ($-1.97$, $-2.00$) represent the
strongest evidence for independent GIPI information content.
\end{tablenotes}
\end{threeparttable}
\end{center}\end{table}

\begin{table}[htbp]\begin{center}
\caption{Check 2: LASSO-QR Surviving Predictors and Pinball Improvement}
\label{tab:lasso}
\begin{threeparttable}\small\setlength{\tabcolsep}{7pt}
\begin{tabular}{@{}cc>{\raggedright\arraybackslash}p{6.2cm}r@{}}
\toprule
$h$ & $\tau$ & Surviving predictors & Impr.(\%)\\
\midrule
\multirow{7}{*}{3}
 & 0.05 & FCI, ip\_yoy                                          & $+3.5$\\
 & 0.10 & FCI, ip\_yoy                                          & $+0.2$\\
 & 0.25 & FCI, ip\_yoy, $\Delta$bkeven$_{5Y}$                   & $+0.3$\\
 & 0.50 & GPR$_z$, ip\_yoy, import\_price\_yoy                  & $+0.8$\\
 & 0.75 & GPR$_z$, ip\_yoy, import\_price\_yoy                  & $+1.5$\\
 & 0.90 & ip\_yoy, import\_price\_yoy                           & $+3.0$\\
 & 0.95 & global\_energy\_yoy                                   & $+8.9$\\
\midrule
\multirow{7}{*}{6}
 & 0.05 & FCI, ip\_yoy                                          & $+7.9$\\
 & 0.10 & FCI, ip\_yoy                                          & $+1.6$\\
 & 0.25 & FCI, ip\_yoy                                          & $+0.1$\\
 & 0.50 & GPR$_z$, ip\_yoy, import\_price\_yoy, food\_yoy       & $+0.2$\\
 & 0.75 & GPR$_z$, ip\_yoy, import\_price\_yoy                  & $+0.1$\\
 & 0.90 & (none survive shrinkage)                              & $+0.2$\\
 & 0.95 & (none survive shrinkage)                              & $+2.2$\\
\midrule
\multirow{7}{*}{12}
 & \cellcolor{yellow!25}0.05 & \cellcolor{yellow!25}\textbf{FCI, import\_price\_yoy}  & \cellcolor{yellow!25}$\mathbf{+13.7}$\\
 & \cellcolor{yellow!25}0.10 & \cellcolor{yellow!25}\textbf{FCI, import\_price\_yoy}  & \cellcolor{yellow!25}$\mathbf{+11.0}$\\
 & 0.25 & FCI, global\_energy\_yoy, $\Delta$bkeven$_{5Y}$       & $+1.3$\\
 & 0.50 & GPR$_z$                                               & $+0.1$\\
 & 0.75 & FCI, ip\_yoy, food\_yoy, $\Delta$bkeven$_{5Y}$        & $+0.1$\\
 & 0.90 & ip\_yoy                                               & $+1.8$\\
 & 0.95 & ip\_yoy                                               & $+5.6$\\
\bottomrule
\end{tabular}
\begin{tablenotes}[flushleft]\footnotesize
\item \textit{Notes:}
$L_1$-penalised quantile regression; 5-fold CV tuning at $\tau=0.10$ per horizon.
Predictor pool: GPR$_z$, FCI, five GIPI components, lagged ip\_yoy.
\textbf{Highlighted}: headline $h=12$ lower-tail results where import prices displace GPR,
consistent with a 9--18-month pass-through lag \citep{conflitti2019}.
\end{tablenotes}
\end{threeparttable}
\end{center}\end{table}

\begin{table}[htbp]\begin{center}
\caption{Check 3: Complete Scoring-Rule Results, Baseline vs.\ Enhanced}
\label{tab:pinball}
\begin{threeparttable}\small\setlength{\tabcolsep}{5.5pt}
\begin{tabular}{@{}cc *{3}{r} @{\hspace{8pt}} *{3}{r}@{}}
\toprule
 &  & \multicolumn{3}{c}{Pinball Loss $\mathcal{L}_\tau$}
    & \multicolumn{3}{c}{Pseudo-$R^2_\tau$}\\
\cmidrule(lr){3-5}\cmidrule(l){6-8}
$h$ & $\tau$ & Base & Enh. & $\Delta$(\%) & Base & Enh. & $\Delta R^2$\\
\midrule
\multirow{7}{*}{3}
 & 0.05 & 0.279 & 0.269 & $+3.5$ & 0.573 & 0.588 & $+0.015$\\
 & 0.10 & 0.399 & 0.398 & $+0.2$ & 0.579 & 0.580 & $+0.001$\\
 & 0.25 & 0.624 & 0.622 & $+0.3$ & 0.545 & 0.546 & $+0.001$\\
 & 0.50 & 0.768 & 0.761 & $+0.8$ & 0.454 & 0.458 & $+0.004$\\
 & 0.75 & 0.623 & 0.614 & $+1.5$ & 0.370 & 0.379 & $+0.009$\\
 & 0.90 & 0.401 & 0.389 & $+3.0$ & 0.272 & 0.293 & $+0.021$\\
 & 0.95 & 0.289 & 0.263 & $+8.9$ & 0.181 & 0.254 & $+0.073$\\
\midrule
\multirow{7}{*}{6}
 & 0.05 & 0.402 & 0.370 & $+7.9$ & 0.384 & 0.433 & $+0.049$\\
 & 0.10 & 0.594 & 0.585 & $+1.6$ & 0.373 & 0.383 & $+0.010$\\
 & 0.25 & 0.956 & 0.955 & $+0.1$ & 0.300 & 0.301 & $+0.001$\\
 & 0.50 & 1.093 & 1.090 & $+0.2$ & 0.217 & 0.219 & $+0.002$\\
 & 0.75 & 0.825 & 0.824 & $+0.1$ & 0.152 & 0.153 & $+0.001$\\
 & 0.90 & 0.504 & 0.503 & $+0.2$ & 0.059 & 0.060 & $+0.001$\\
 & 0.95 & 0.334 & 0.327 & $+2.2$ & 0.029 & 0.050 & $+0.021$\\
\midrule
\multirow{7}{*}{12}
 & \cellcolor{yellow!25}0.05 & \cellcolor{yellow!25}0.559 & \cellcolor{yellow!25}0.482 & \cellcolor{yellow!25}$\mathbf{+13.7}$ & \cellcolor{yellow!25}0.146 & \cellcolor{yellow!25}0.263 & \cellcolor{yellow!25}$\mathbf{+0.117}$\\
 & \cellcolor{yellow!25}0.10 & \cellcolor{yellow!25}0.887 & \cellcolor{yellow!25}0.790 & \cellcolor{yellow!25}$\mathbf{+11.0}$ & \cellcolor{yellow!25}0.066 & \cellcolor{yellow!25}0.168 & \cellcolor{yellow!25}$\mathbf{+0.102}$\\
 & 0.25 & 1.347 & 1.330 & $+1.3$ & 0.014 & 0.027 & $+0.013$\\
 & 0.50 & 1.358 & 1.357 & $+0.1$ & 0.027 & 0.028 & $+0.001$\\
 & 0.75 & 0.944 & 0.943 & $+0.1$ & 0.026 & 0.027 & $+0.001$\\
 & 0.90 & 0.492 & 0.483 & $+1.8$ & 0.084 & 0.101 & $+0.017$\\
 & 0.95 & 0.291 & 0.275 & $+5.6$ & 0.162 & 0.209 & $+0.047$\\
\bottomrule
\end{tabular}
\begin{tablenotes}[flushleft]\footnotesize
\item \textit{Notes:} Highlighted rows = headline $(h{=}12,\ \tau{\leq}0.10)$ results.
Enhanced weakly dominates at every cell; minimum improvement $+0.08\%$.
Improvements concentrated in tails; median ($\tau{=}0.50$) improvements near zero at all horizons.
In-sample, 1997M01--2025M12.
\end{tablenotes}
\end{threeparttable}
\end{center}\end{table}

\FloatBarrier

% ═══════════════════════════════════════════════════════════════════════════════
\section{Layer 0: Real-Time Event Detection}\label{sec:layer0}
% ═══════════════════════════════════════════════════════════════════════════════

\subsection{GDELT 2.0 News Surveillance}

The Global Database of Events, Language, and Tone (GDELT 2.0,
\citealp{leetaru2013}) monitors broadcast, print, and online news
in over 100 languages, updated every 15 minutes.
We query the GDELT Doc API~v2 with a Boolean AND of Middle East location
terms (Iran, Israel, Gaza, Lebanon, Yemen, Strait of Hormuz, Red Sea,
Persian Gulf, Saudi Arabia) crossed with conflict vocabulary
(military, strike, attack, missile, bomb, naval, escalation, drone).
The lookback window is 48 hours in calm regime, extended to 168 hours
in elevated or crisis regime.
Up to 250 articles are returned, sorted by recency.
Each article is scored for location relevance and CAMEO conflict severity;
the top-20 by combined score are passed to the LLM classifier.

\subsection{LLM CAMEO Coding}

The 20 ranked headlines are sent to an LLM (Anthropic Claude) with the
CAMEO severity guide \citep{schrodt2009}, the current GPR $z$-score as
an anchor, and instructions to return validated JSON with fields:
\texttt{severity}~(0--10), \texttt{cameo\_codes}~(list),
\texttt{gpr\_z\_estimate}, \texttt{deescalation}~(bool),
and \texttt{summary} ($\leq 100$ words).
A regex fallback operates on API unavailability.
The composite severity score is:
\begin{equation}
S_t = 0.5\,s^{\text{LLM}}_t + 0.3\,s^{\text{rule}}_{t,P_{75}}
    + 0.2\,\mathbf{1}[\text{AIS anomaly}_t].
\label{eq:severity}
\end{equation}

\subsection{AIS Tanker Equity Proxy}

We proxy Strait of Hormuz disruption using four listed VLCC operators
(INSW, TK, TRMD, FRO) and the Brent--WTI spread.
The anomaly flag fires when both the 30-day rolling $z$-score of the
tanker basket \textit{and} the Brent--WTI spread $z$-score simultaneously
exceed $1.5\sigma$.

\subsection{Bridging the Publication Lag}
\label{sec:lag}

The GPR index is published with a one-month lag: as of March 2026, the
most recent official release covers November 2025.
This creates a blind spot for real-time risk monitoring---precisely the
period in which fast-moving geopolitical developments are most
policy-relevant.

Layer~0 bridges this gap by producing a \textit{sub-monthly GPR nowcast}
$\hat{z}_t$ derived from live GDELT headlines, LLM CAMEO severity scoring,
and the AIS tanker proxy.
The nowcast is injected into the GaR model as the current-month
$\text{GPR}_z$ observation, replacing the stale monthly reading with a
contemporaneous estimate.
Formally, define the publication lag as $\ell = 1$ month.
The nowcast augments the information set from
$\mathcal{I}_{t-\ell}$ (monthly GPR only) to
$\mathcal{I}_t$ (monthly GPR + real-time Layer~0 signal),
so the enhanced GaR forecast conditions on
\begin{equation}
\hat{z}_t \;=\; f\!\left(\,S_t,\;\text{AIS}_t,\;\hat{z}^{\text{GPR}}_{t-\ell}\right),
\label{eq:nowcast}
\end{equation}
where $S_t$ is the composite severity score (Equation~\ref{eq:severity}),
$\text{AIS}_t$ is the tanker equity proxy, and
$\hat{z}^{\text{GPR}}_{t-\ell}$ is the last official GPR $z$-score.

As of March 2026, the official GPR reads $z=2.65$ (November 2025,
\textsc{crisis} band), while the Layer~0 nowcast yields $\hat{z}=1.97$
(\textsc{elevated} band).
The gap of $0.68\sigma$ captures within-month de-escalation dynamics
invisible to the monthly release schedule: military posturing and movement
restrictions are ongoing (CAMEO codes 15/191) but no kinetic chokepoint
attack has materialised and the AIS dual trigger remains inactive.
Without Layer~0, the GaR model would overstate tail risk by anchoring on
the stale crisis reading; with it, the six-month recession probability
falls from 16.7\% (baseline, unadjusted) to 2.6\% (enhanced, Layer~0
nowcast injected).

\subsection{Current Nowcast}

Figure~\ref{fig:layer0} summarises the full Layer~0 output panel as of
March 7, 2026.

\begin{figure}[htbp]\centering
\includegraphics[width=\textwidth]{fig09_layer0.pdf}
\caption{Layer~0 Real-Time Event Detection: March 7, 2026.
\textit{Top-left}: GPR $z$-score time series (2023--2025) with regime shading
(green = calm, amber = elevated, red = crisis); horizontal dotted lines at
elevated ($z=1.5$) and crisis ($z=2.5$) thresholds.
November 2025 annotated: $z=2.65$.
\textit{Top-right}: CAMEO severity scale used by the LLM classifier,
with a vertical red dashed line at the current composite severity of 5.6/10.
Active codes 15 (military posture) and 191 (blockade/restrict movement) are
highlighted; note the gap below the aerial-weapons threshold (8.0) confirming
no kinetic attack has occurred.
\textit{Bottom-left}: Arab Light--WTI spread over time with historical mean
(dashed). No anomalous widening visible; the spread remains near its average
of approximately $-\$46$, confirming no Hormuz premium.
\textit{Bottom-right}: Nowcast summary table with all key metrics. GIPI
reading of $-1.28$ (CALM) annotated alongside the ELEVATED GPR regime,
documenting the decoupling in real-time.}
\label{fig:layer0}
\end{figure}

Table~\ref{tab:l0} presents the March 7, 2026 signal numerically.

\begin{table}[htbp]
\caption{Layer~0 Real-Time Signal, March 7, 2026}
\label{tab:l0}
\begin{center}
\begin{threeparttable}\small\setlength{\tabcolsep}{6pt}
\begin{tabular}{@{}p{4.2cm} p{3.4cm} p{5.8cm}@{}}
\toprule
\textit{Metric} & \textit{Value} & \textit{Interpretation} \\
\midrule
\multicolumn{3}{l}{\textit{A. Real-time event signal}}\\[2pt]
Regime & \textsc{elevated} & GPR nowcast $z\in[1.5,\,2.5)$\\
Composite severity $S_t$ & 5.6\,/\,10 & Military posture; no kinetic strike\\
GPR nowcast $\hat{z}$ & 1.97 & Near upper bound of elevated band\\
CAMEO codes (dominant) & 15, 191 & Force buildup; movement restriction\\
GDELT articles & 3 (100\% ME) & All articles Middle East--relevant\\
AIS anomaly flag & \textbf{No} & Both $z$-scores $< 1.5\sigma$\\
Tanker basket $z$ & $-0.97$ & Below anomaly threshold\\
Brent--WTI spread & \$1.8 & No Hormuz supply premium\\
\midrule
\multicolumn{3}{l}{\textit{B. Monthly anchors (most recent releases)}}\\[2pt]
GPR level (Nov 2025) & 371.1, $z=2.65$ & \textsc{crisis} band; 1-month lag\\
GIPI (Dec 2025) & $-1.28$ & \textsc{calm}: all channels subdued\\
\midrule
\multicolumn{3}{l}{\textit{C. Implied 6-month GaR nowcast (enhanced model)}}\\[2pt]
Median IP YoY & $+0.9\%$ & Benign central case\\
GaR$_5$ (5th percentile) & $-3.5\%$ & Left tail remains material\\
$\Pr(\text{recession})$ & 2.6\% & Low but not negligible\\
Conditional skewness & $-0.21$ & Channel re-engagement risk priced\\
\bottomrule
\end{tabular}
\begin{tablenotes}[flushleft]\footnotesize
\item \textit{Notes:}
CAMEO code 15 = exhibit military posture or force buildup;
code 191 = impose blockade or restrict movement.
Neither aerial weapons (195) nor mass violence (200) codes are active,
confirming no kinetic chokepoint attack has occurred.
The gap between monthly GPR ($z=2.65$, November~2025) and Layer~0 nowcast
($\hat{z}=1.97$) captures within-month de-escalation dynamics invisible
to the monthly release schedule.
AIS dual trigger: 30-day rolling $z > 1.5$ required for \textit{both}
the VLCC equity basket return \textit{and} the Brent--WTI spread; neither
condition is met.
Layer~0 nowcast $\hat{z}=1.97$ is injected as the current-month GPR observation
in the GaR model, replacing the stale November 2025 official release ($z=2.65$)
and reducing the implied six-month recession probability from 16.7\% to 2.6\%.
\end{tablenotes}
\end{threeparttable}
\end{center}
\end{table}

The CAMEO code combination 15/191 (military posture + blockade) with
no AIS anomaly and Brent--WTI spread at \$1.8 is the diagnostic signature
of a \textit{coercive posturing} episode: military deployments and
movement restrictions are ongoing but no physical chokepoint disruption
has occurred.
A genuine Hormuz disruption would register CAMEO 195 (aerial weapons)
or 200, activate the AIS dual trigger, and widen the Brent--WTI spread
sharply---as occurred briefly after the September 2019 Abqaiq attack.
None of those signals are present on March 7, 2026.

% ═══════════════════════════════════════════════════════════════════════════════
\section{Conclusion}\label{sec:conclude}
% ═══════════════════════════════════════════════════════════════════════════════

\noindent\textbf{What this paper does.}
We introduce GIPI, a monthly index of geopolitical inflation transmission
pressure---the first principal component of five channel variables---and
integrate it into a Growth-at-Risk framework via a GPR$\times$GIPI
interaction term.
The core idea is simple: a geopolitical shock is most dangerous to the
growth distribution when the inflation channels are open (high GIPI), and
least dangerous when they are closed (low GIPI, as in 2025).

\noindent\textbf{What we find.}
Four results emerge.
\textit{First}, the amplification mechanism is statistically robust:
GIPI carries independent left-tail information in every one of 21
quantile-horizon robustness cells.
The orthogonalised coefficient at $h=12$, $\tau=0.05$ is $-1.97$
($p<0.01$)---the single strongest result in the paper.
\textit{Second}, augmenting the \citet{adrian2019} GaR model with GIPI
reduces 12-month tail pinball loss by 13.7\% and nearly doubles
pseudo-$R^2$ at the fifth quantile, with LASSO identifying import
price growth as the dominant long-horizon channel.
\textit{Third}, the 2023--25 GPR--GIPI decoupling ($\rho=0.006$) is
unprecedented in the post-2000 sample and is the reason why crisis-level
geopolitical risk has not triggered a recession: US shale production,
IEA reserve coordination, and Fed credibility have simultaneously
suppressed all five transmission channels.
The model correctly computes a 2.6\% six-month recession probability for
this configuration, against 16.7\% from the na\"ive baseline.
\textit{Fourth}, Layer~0 confirms the decoupling persists in real time:
CAMEO codes 15/191 (posturing, no kinetics), AIS dual-trigger inactive,
Brent--WTI at \$1.8---none of the signatures of a genuine chokepoint
disruption are present.

\noindent\textbf{When to worry.}
The $-0.21$ conditional skewness at $h=6$ is the paper's key risk-management
statistic: the left tail remains meaningfully fat even when the median is
benign.
The decoupling breaks down if \textit{any} of three conditions changes:
(i)~US production falls below 12\,mb/d (possibly from tariff-driven capex
cuts); (ii)~IEA reserve buffers are depleted; or (iii)~inflation expectations
de-anchor, as measured by a GIPI reading above $+2$.
The Layer~0 pipeline will detect (ii) and (iii) within 48 hours via the AIS
anomaly trigger and the breakeven component of GIPI respectively.

\noindent\textbf{Extensions.}
Three planned extensions would strengthen the framework.
Real vessel-level AIS data (Kpler or Spire Maritime) would replace the
equity proxy for a sharper Hormuz signal.
Prediction market probabilities (Polymarket, Metaculus) could weight
Layer~0 severity scores by calibrated escalation-to-kinetics probabilities.
State-dependent local projections---estimated separately in high-GIPI
and low-GIPI regimes---would formally test whether the amplification
hypothesis operates symmetrically at the IRF level, complementing the
quantile GaR evidence presented here.

\newpage

% ── References ────────────────────────────────────────────────────────────────
\bibliographystyle{aer}
\begin{thebibliography}{99}

\bibitem[Adrian et al.(2019)]{adrian2019}
Adrian, T., N.~Boyarchenko, and D.~Giannone. 2019.
``Vulnerable Growth.'' \textit{American Economic Review}
109(4): 1263--1289.

\bibitem[Baumeister and Kilian(2016)]{baumeister2016}
Baumeister, C., and L.~Kilian. 2016.
``Forty Years of Oil Price Fluctuations.''
\textit{Journal of Economic Perspectives} 30(1): 139--160.

\bibitem[Benigno et al.(2022)]{benigno2022}
Benigno, G., J.~di Giovanni, J.J.J.~Groen, and A.I.~Noble. 2022.
``The GSCPI: A New Barometer of Global Supply Chain Pressures.''
FRBNY Staff Reports No.\ 1017.

\bibitem[Bloom(2009)]{bloom2009}
Bloom, N. 2009. ``The Impact of Uncertainty Shocks.''
\textit{Econometrica} 77(3): 623--685.

\bibitem[Bonadio et al.(2021)]{bonadio2021}
Bonadio, B., Z.~Huo, A.A.~Levchenko, and N.~Pandalai-Nayar. 2021.
``Global Supply Chains in the Pandemic.''
\textit{Journal of International Economics} 133: 103534.

\bibitem[Caldara and Iacoviello(2022)]{caldara2022}
Caldara, D., and M.~Iacoviello. 2022.
``Measuring Geopolitical Risk.''
\textit{American Economic Review} 112(4): 1194--1225.

\bibitem[Caldara et al.(2025)]{caldara2025}
Caldara, D., S.~Conlisk, M.~Iacoviello, and M.~Penn. 2025.
``Do Geopolitical Risks Raise or Lower Inflation?''
\textit{Journal of Monetary Economics}, forthcoming.

\bibitem[Carri\`{e}re-Swallow et al.(2023)]{carriereswallow2023}
Carri\`{e}re-Swallow, Y., N.C.~Firat, D.~Furceri, and D.~Jimenez. 2023.
``State-Dependent Exchange Rate Pass-Through.''
\textit{Journal of International Money and Finance} 134: 102830.

\bibitem[Jord\`{a}(2005)]{jorda2005}
Jord\`{a}, \`{O}. 2005.
``Estimation and Inference of Impulse Responses by Local Projections.''
\textit{American Economic Review} 95(1): 161--182.

\bibitem[Kilian(2009)]{kilian2009}
Kilian, L. 2009.
``Not All Oil Price Shocks Are Alike.''
\textit{American Economic Review} 99(3): 1053--1069.

\bibitem[Koenker and Bassett(1978)]{koenker1978}
Koenker, R., and G.~Bassett. 1978.
``Regression Quantiles.'' \textit{Econometrica} 46(1): 33--50.

\bibitem[Koenker and Machado(1999)]{koenker1999}
Koenker, R., and J.A.F.~Machado. 1999.
``Goodness of Fit and Related Inference Processes for Quantile Regression.''
\textit{Journal of the American Statistical Association}
94(448): 1296--1310.

\bibitem[Leetaru and Schrodt(2013)]{leetaru2013}
Leetaru, K., and P.A.~Schrodt. 2013.
``GDELT: Global Data on Events, Location, and Tone, 1979--2012.''
\textit{ISA Annual Conference}. San Francisco.

\bibitem[L\'{o}pez-Salido and Loria(2024)]{lopezloria2024}
L\'{o}pez-Salido, D., and F.~Loria. 2024.
``Inflation at Risk.''
\textit{Journal of Monetary Economics} 145: 103568.

\bibitem[Newey and West(1987)]{newey1987}
Newey, W.K., and K.D.~West. 1987.
``A Simple, Positive Semi-definite, Heteroskedasticity and
Autocorrelation Consistent Covariance Matrix.''
\textit{Econometrica} 55(3): 703--708.

\bibitem[Schrodt(2009)]{schrodt2009}
Schrodt, P.A. 2009.
``CAMEO: Conflict and Mediation Event Observations Codebook.''
Version 1.1b3. Pennsylvania State University.

\bibitem[Conflitti and Luciani(2019)]{conflitti2019}
Conflitti, C., and M.~Luciani. 2019.
``Oil Price Pass-Through into Core Inflation.''
\textit{The Energy Journal} 40(6): 221--248.


\bibitem[Adams et al.(2021)]{adams2021}
Adams, P.A., T.~Adrian, N.~Boyarchenko, and D.~Giannone. 2021.
``Forecasting Macroeconomic Risks.''
\textit{International Journal of Forecasting} 37(3): 1173--1191.

\bibitem[Aiyar et al.(2023)]{aiyar2023}
Aiyar, S., A.~Ilyina, and others. 2023.
``Geoeconomic Fragmentation and the Future of Multilateralism.''
IMF Staff Discussion Note SDN/2023/001.

\bibitem[Baker, Bloom and Davis(2016)]{baker2016}
Baker, S.R., N.~Bloom, and S.J.~Davis. 2016.
``Measuring Economic Policy Uncertainty.''
\textit{Quarterly Journal of Economics} 131(4): 1593--1636.

\bibitem[Baumeister and Hamilton(2019)]{baumeister2019}
Baumeister, C., and J.D.~Hamilton. 2019.
``Structural Interpretation of Vector Autoregressions with Incomplete
Identification: Revisiting the Role of Oil Supply and Demand Shocks.''
\textit{American Economic Review} 109(5): 1873--1910.

\bibitem[Caldara, Iacoviello and Yu(2025)]{caldara2025shortages}
Caldara, D., M.~Iacoviello, and D.~Yu. 2025.
``Measuring Shortages since 1900.''
International Finance Discussion Papers No.~1407,
Board of Governors of the Federal Reserve System.

\bibitem[Chavleishvili and Manganelli(2024)]{chavleishvili2024}
Chavleishvili, S., and S.~Manganelli. 2024.
``Forecasting and Stress Testing with Quantile Vector Autoregression.''
\textit{Journal of Applied Econometrics} 39(1): 66--85.

\bibitem[Federle et al.(2024)]{federle2024}
Federle, J., A.~Meier, G.J.~M\"{u}ller, W.~Mutschler, and M.~Schularick. 2024.
``The Price of War.''
CEPR Discussion Paper No.~18834 / Kiel Working Paper No.~2262.

\bibitem[Fern\'{a}ndez-Villaverde, Mineyama and Song(2024)]{fernandez2024}
Fern\'{a}ndez-Villaverde, J., T.~Mineyama, and D.~Song. 2024.
``Are We Fragmented Yet? Measuring Geopolitical Fragmentation and Its
Causal Effects.'' NBER Working Paper No.~32638.

\bibitem[Hassan et al.(2019)]{hassan2019}
Hassan, T.A., S.~Hollander, L.~van~Lent, and A.~Tahoun. 2019.
``Firm-Level Political Risk: Measurement and Effects.''
\textit{Quarterly Journal of Economics} 134(4): 2135--2202.

\bibitem[Metiu(2025)]{metiu2025}
Metiu, N. 2025.
``Global Financial Transmission of Geopolitical Risk Shocks.''
Deutsche Bundesbank / SSRN Working Paper No.~5195860.

\bibitem[Pinchetti(2024)]{pinchetti2024}
Pinchetti, M. 2024.
``Geopolitical Risk and Inflation: The Role of Energy Markets.''
Centre for Macroeconomics Discussion Paper No.~2431.
(Also Banque de France Working Paper No.~1005, 2025.)

\bibitem[Plagborg-M{\o}ller and Wolf(2021)]{plagborg2021}
Plagborg-M{\o}ller, M., and C.K.~Wolf. 2021.
``Local Projections and VARs Estimate the Same Impulse Responses.''
\textit{Econometrica} 89(2): 955--980.

\end{thebibliography}

\end{document}
